% ============================================================
% 5_E_admin_teacher.tex — 管理员端 + 教师端功能需求
% 编写人：E  日期：2026-07-15
% 职责模块：管理员后台管理 / 教师批阅管理 / 账号管理
% ============================================================

\section{管理员端功能需求}

\subsection{管理员端功能模块划分}

管理员端承担系统基础数据的维护和论文模板的配置管理职责，是系统的后台管理核心。管理员端划分为以下几个子模块：

\begin{enumerate}
    \item \textbf{学院信息管理}：对学校各学院信息进行增删改查，学院信息是用户注册和模板绑定的基础；
    \item \textbf{用户账号管理}：查看系统全部用户列表，支持禁用/启用用户账号，查询用户注册信息；
    \item \textbf{模板管理}：创建、编辑、启用/停用论文模板，管理模板版本；
    \item \textbf{系统数据概览}：查看系统运行数据总览（用户数量、论文数量、模板数量等）。
\end{enumerate}

\subsection{学院信息管理}

\reqitem{FR-E01}{管理员可创建新的学院信息，填写学院名称（如"计算机科学与技术学院"）和学院代码（如"CS"，全局唯一）。学院代码一经创建不可修改，学院名称可编辑。}
\reqitem{FR-E02}{管理员可查看全部学院列表，支持按学院名称搜索。列表展示学院名称、学院代码、关联模板数量、关联用户数量。}
\reqitem{FR-E03}{管理员可编辑已有学院名称，或删除无关联数据的学院。删除前系统自动校验该学院是否关联了模板或用户，若存在关联则阻止删除并提示关联数量。}

\subsection{用户账号管理}

\reqitem{FR-E04}{管理员可查看系统全部用户列表，支持按角色（ADMIN/STUDENT/TEACHER）和学院进行筛选。列表展示用户名、真实姓名、角色、所属学院、注册时间、账号状态。}
\reqitem{FR-E05}{管理员可将异常账号（如恶意注册、长期未使用的账号）禁用。禁用后该账号无法登录系统，已登录的Token在Redis中的缓存被清除。支持重新启用已禁用的账号。}
\reqitem{FR-E06}{管理员可重置指定学生的密码为初始密码（123456），重置后该学生下次登录系统时会收到修改密码的提示。}

\subsection{系统数据概览}

\reqitem{FR-E07}{管理员登录后进入系统数据概览仪表盘页面，展示以下统计数据：系统注册用户总数（按角色分类统计）、论文模板总数（按类型分类统计）、在编论文总数（按状态分类统计）、当日导出任务数。数据以卡片和图表形式展示，支持按时间范围筛选。}

% ============================================================
\section{教师端功能需求}

\subsection{教师端功能模块划分}

教师端是面向论文批阅者的操作入口，承载论文审阅、批注添加、评分与退回等核心功能。教师端划分为以下几个子模块：

\begin{enumerate}
    \item \textbf{批阅任务管理}：查看待批阅论文列表、进行中批阅、已完成批阅历史；
    \item \textbf{论文在线批阅}：在线查看论文全文内容，在指定文字位置添加批注，查看已有批注列表；
    \item \textbf{评分与退回}：对论文进行综合评分和等级评定，或退回论文要求学生修改重提；
    \item \textbf{批阅历史查询}：查看历次批阅记录、分数变化轨迹、批注修改对比。
\end{enumerate}

\subsection{批阅任务管理}

\reqitem{FR-E08}{教师登录后进入批阅任务列表页，展示已提交至该教师的全部论文。列表支持按学院筛选、按提交时间排序，展示字段包括：论文标题、学生姓名、所属学院、提交时间、当前状态（待批阅/已批阅/已退回）。}
\reqitem{FR-E09}{教师可点击论文进入批阅详情页。支持在同一页面中顺序浏览多个批阅任务（上一个/下一个），提升批量批阅效率。}

\subsection{论文在线批阅}

\reqitem{FR-E10}{教师进入批阅页面后，左侧展示论文的章节树，右侧展示论文全文HTML渲染内容（只读模式，不可编辑）。教师可通过点击章节树切换查看各章节内容。}
\reqitem{FR-E11}{教师可在论文正文中选中任意文字段落，选中的文字以高亮背景标记，同时弹出批注输入框。教师填写批注内容后确认提交，该批注以卡片形式显示在右侧批注列表中。}
\reqitem{FR-E12}{批注以锚点形式与原文位置绑定，教师点击批注列表中的某条批注，页面自动滚动并高亮对应的原文位置。批注支持编辑和删除操作。}

\subsection{评分与退回}

\reqitem{FR-E13}{教师完成全部批注后，点击"评阅完成"按钮弹出评分对话框。评分要素包括：整体评语（可选填，支持HTML格式）、打分（0-100整数滑块）、等级（系统自动换算）。}
\reqitem{FR-E14}{教师可以选择"通过"（REVIEWED）或"退回"（RETURNED）两种操作。选择退回时必须填写退回原因，且字数不少于10个字符。退回后论文状态变为RETURNED，is\_locked改为FALSE，学生可重新编辑。}
\reqitem{FR-E15}{论文通过后状态变为REVIEWED，学生可查看教师的评语、得分和等级，以及全部批注的详细内容。}}

\subsection{批阅历史查询}

\reqitem{FR-E16}{教师可查看某篇论文的历次批阅记录列表，包括每次批阅的分数、等级、操作类型（通过/退回）、操作时间和对应的退回原因。支持两次批阅的得分变化对比展示。}
\reqitem{FR-E17}{学生可查看历次批阅的全部批注内容和评语，系统在界面上清晰标注哪些批注是本次新增的、哪些是已读的。
